\section{The direct zero-entrance near equation}

Use the normalized cylinder spaces
\[
 H=L^2(s^{-3/2}\,ds\,dq/c\,d\eta),\qquad
 H_T=L^2(s^{-1/2}\,ds\,dq/c\,d\eta),\qquad
 Z=Z_1\oplus\rho_b^{-1/2}H_T,
\]
where \(Z_1=L^2(s^{-1/2}W^2\,ds\,dq/c\,d\eta)\) and
\(W=\langle\eta/\sqrt{1+s}\rangle^{m_Z}\), \(m_Z>3/2\).
The stress inverse always acts on the sum of the components. Let
\(\mathcal M_0\) be multiplication by the actual flat train, with
zero second output component, and let \(\mathcal M_U\) be the
exact thin/repair multiplier. Set
\[
 \mathsf D_0=I+\mathcal M_0\mathcal S_0,\qquad
 g(s)=e^{\gamma s/(1+\nu s)},\qquad \mathsf G T=gT.
\]
The direct causal construction on the full moving clock gives
\begin{equation}\label{near:flat-bound}
 \|\mathsf G^{-1}\mathsf D_0^{-1}\mathsf G\|_{Z\to Z}
 \le Q:=P_\nu.
\end{equation}
There are no auxiliary entrance coordinates in this statement.
The ordinary norm of this inverse is not asserted to have the same
bound; it can contain the full planar amplification. The polynomial
degree is fixed by the velocity and tensor spaces.

The finite-curvature free inverse is the conjugate of the Cartesian
weak-stress inverse. In particular its divergence and pressure are
not replaced by planar ones. In the profile chart
\[
 s=R^2/\rho^2,\quad q=R\log(\rho/R)+R\tau/2,
 \quad\eta=Rz/\rho,\quad R=\epsilon^{-1},
\]
its exact divergence is
\[
 \operatorname{div}(\sqrt{s}\,b)
 =s\{[\partial_q-\epsilon(2s\partial_s+\eta\partial_\eta)]b_r
                   +\partial_\eta b_z\}.
\]
The geometric comparison below concerns these exact operators.

Fix a small positive logarithmic gap \(\delta\). For an output age
\(s\) and input age \(t\), define
\[
 \mathcal S_{\rm far}
   =1_{\{s<S,\ t>e^\delta s\}}\mathcal S_\epsilon,
 \quad \mathcal S_N=\mathcal S_\epsilon-\mathcal S_{\rm far},
 \quad \mathsf D_N=I+\mathcal M_U\mathcal S_N.
\]
The notation for the mask means its exact kernel realization with
tensor partitions before distributional divergence. Interface terms
are retained. A masked intermediate velocity need not satisfy the
whole-space divergence-free constraint by itself; the stress equation is
well defined independently of that issue. Only the recombined full
Cartesian inverse is used to recover the final physical mode.

This \(\mathsf D_N\) is distinct from subtracting the divergence-free
old-to-young velocity in the nilpotent translation transformation.
Their difference is retained in the transformed equation.

\begin{proposition}\label{near:equation-comparison}
Let \(\mu\) range over a fixed compact subset of \(0<\Re\mu<1/4\).
Put \(s_-=\epsilon^a\), \(S=\epsilon^{-b}\), \(0<b<1\).
There is a fixed integer \(K\), chosen before \(a,b\), such that
\begin{equation}\label{near:deltaN}
 \begin{split}
 \delta_N:=\|\mathsf G^{-1}(\mathsf D_N-\mathsf D_0)\mathsf G\|
 \le P_\nu\big[&\epsilon^{a/2}+\epsilon^{b/2}
       +\epsilon e^{(m_0+\eta_\delta)/\nu}
       +\epsilon^{1/50-K(a+b)}e^{\eta_\delta/\nu}\\
       &+\epsilon^{3/22}e^{(\eta_\delta+m_0/2)/\nu}
       +e^{-(D_\perp-\eta_\delta)/\nu}\big],
 \end{split}
\end{equation}
where \(\eta_\delta\downarrow0\) as \(\delta\downarrow0\).
The constant \(D_\perp>0\) specifies the transverse-tail accuracy.
If \(Q\delta_N<1/2\), then
\begin{equation}\label{near:inverse-comparison}
 \begin{split}
 \|\mathsf G^{-1}\mathsf D_N^{-1}\mathsf G\|&\le2Q,\\
 \|\mathsf G^{-1}(\mathsf D_N^{-1}-\mathsf D_0^{-1})\mathsf G\|
       &\le2Q^2\delta_N.
 \end{split}
\end{equation}
Both statements are equation-only estimates; no terminal functional
or affine source is added.
\end{proposition}

\begin{proof}
The required free comparison is global in its input and compact in
its output:
\[
 \|\chi(\mathcal S_\epsilon-\mathcal S_0)\|_{H_T\to H}
       \le C\Gamma^K\epsilon^{1/50}.
\]
Here \(\Gamma\) is a polynomial in inverse lower age, upper age,
transverse box size, viscosity inverse, cell width, and finitely many
cutoff derivatives. The exponent \(K\) is fixed. The transverse
cutoff may be chosen at \(L_\perp=C_\perp(c/\nu+\sqrt S+1)\),
so \(L_\perp\le C'_\perp S\) for small \(\nu\), and
\(\Gamma^K\le P_\nu\epsilon^{-K(a+b)}\).
The actual flat train coefficient outside this box is at most
\(P_\nu e^{-D_\perp/\nu}\). Its curved stream and entrance
adjustments are estimated separately by their global small multiplier
norm, so their tails need not be Gaussian. No transverse weight is
imposed on the broad repair coefficient.

Partition \([s_-,S]\) into output bins \([l,r]\) with
\(r/l<e^{\delta/2}\). Their number is a viscosity polynomial
in the distinguished regime. The retained input region on a bin
is the complete block \(t<e^\delta l\) and the partial strip
\(e^\delta l<t<e^\delta r\). On the complete block, multiplication
by \(g\) on inputs and \(g^{-1}\) on outputs costs at most
\(g(e^\delta l)/g(l)\). More generally
\[
 \sup_{s>0,\ s\le t\le Ms}
 \nu\left(\frac{t}{1+\nu t}-\frac{s}{1+\nu s}\right)
     =\frac{\sqrt M-1}{\sqrt M+1}.
\]
Indeed, writing \(t=ks\), \(x=\nu s\), the expression is
\((k-1)x/[(1+x)(1+kx)]\), maximized at \(x=k^{-1/2}\).
Thus all complete blocks and partial strips cost at most
\(e^{\eta_\delta/\nu}\), with
\(\eta_\delta=\gamma\tanh(3\delta/8)\) after enlarging the
overlapping bins if necessary. The polynomial bin count is included
in \(P_\nu\). Applying the compact-output comparison and the train
multiplier proves the fourth term of \eqref{near:deltaN}; the
transverse tail gives its last term.

On a partial strip the flat causal inverse is zero. Its source radius
is separated from the output radius by a fixed multiple of
\(R/\sqrt r\), also after the source contraction in the resolvent.
By \eqref{sep:kernel}, the heat--Leray stress kernel has the
uniform positive majorant $C|x|^{-4}$, whose spatial $L^1$ norm
outside distance $d$ is $C/d$.
The exact dilation integral on an \(L^2\) source has power
\(\lambda^{2\Re\mu+1/2}\), hence a uniform integrability margin.
Young's inequality, followed by normalized-phase Minkowski, gives
\(P_\nu\epsilon\sqrt r\) for this output-dependent strip.
After gain multiplication this is at most
\(P_\nu e^{\eta_\delta/\nu}\epsilon^{1-b/2}\), absorbed by
the fourth term. This direct positive-kernel estimate is needed;
boundedness of an arbitrary triangular operator truncation is not used.

For \(s<s_-\), the train multiplier has norm \(P_\nu\sqrt{s_-}\).
Retained inputs have \(t<e^\delta s_-\), where the gain ratio
is \(1+o(1)\). Subtracting the forbidden separated subblock costs
at most \(P_\nu\epsilon\sqrt{s_-}\), by the same positive-kernel
estimate. For \(s>S\), use the full free inverse and the train tail
\(P_\nu(S^{-1/2}+\epsilon)\); its gain ratio is bounded by
\[
 \exp\{\gamma/[\nu(1+\nu S)]\}=1+o(1).
\]
These give the first two terms.

The same complete blocks, using free energy in place of the difference
estimate, give a global bound without a coefficient multiplier. The
positive partial-strip bounds remain valid, the lower end has gain
ratio \(1+o(1)\), and the upper end has the ratio just displayed.
Consequently
\begin{equation}\label{near:masked-free-bound}
 \|g^{-1}\mathcal S_N g\|_{H_T\to H}
       \le P_\nu e^{\eta_\delta/\nu}.
\end{equation}

The exact thin coefficient differs from the flat train by its
\(\epsilon\) stream completion and fixed entrance adjustment,
whose multiplier norm is \(P_\nu\epsilon e^{m_0/\nu}\).
For its contribution use the exact decomposition
\[
 \mathcal M_{\rm thin}\mathcal S_N-\mathcal M_0\mathcal S_0
 = (\mathcal M_{\rm thin}-\mathcal M_0)\mathcal S_N
       +\mathcal M_0(\mathcal S_N-\mathcal S_0).
\]
The first uses the global near free bound and costs
\(P_\nu\epsilon e^{(m_0+\eta_\delta)/\nu}\), including
all tails of the entrance adjustment. The second is the geometric
difference with the Gaussian train multiplier estimated above.
Each broad input/output coupling is bounded by a polynomial times
\(\sqrt{\rho_b}\), using the unweighted free stress bound.
The same bin and end-region estimates apply without a transverse
weight; since \(\rho_b=P_\nu\epsilon^{3/11}e^{m_0/\nu}\),
this gives the fifth term. This proves the complete equation
difference, including every split component.

Conjugate by \(\mathsf G\) and apply the geometric inverse series
to \(\mathsf D_N=\mathsf D_0+(\mathsf D_N-\mathsf D_0)\).
Its norm is at most \(Q\delta_N<1/2\).
The resolvent identity then gives \eqref{near:inverse-comparison}.
\end{proof}

The parameters can be chosen without a terminal observation. Fix the
derivative budgets and \(K\), then choose \(a,b>0\) with
\(K(a+b)<1/100\). Let
\[
 \theta=\min\{a/2,b/2,1/100,3/22\},\qquad d_0=\gamma/3.
\]
Choose \(\delta,m_0\) so that
\(\eta_\delta<\theta d_0/8\),
\(m_0<\theta d_0/8\), and take
\(D_\perp>\eta_\delta+\theta d_0\).
The background construction realizes the prescribed \(m_0\).
All these choices precede \(\nu\downarrow0\).
For \(\epsilon=C_*\nu^r e^{-d_0/\nu}\), uniformly on bounded
\(r\) and positive compact \(C_*\)-ranges,
\[
 \delta_N\le P_\nu e^{-\theta d_0/(2\nu)},\qquad
 Q^2\delta_N\le P_\nu e^{-\theta d_0/(2\nu)}.
\]
Thus the inverse and its comparison close with a positive exponential
margin. The choice is uniform over the spectral interval while all
geometric parameters and the exact background are held fixed.

\subsection{The actual forced column}

Let \(f\) be a normalized body source, so that its physical
representative is \(\mathcal N_v f\), and set
\[
 u_\epsilon=\mathcal M_U\mathcal R_\epsilon f,\qquad
 T_\epsilon=\mathsf D_N^{-1}u_\epsilon,
 \qquad b_\epsilon=\mathcal R_\epsilon f-\mathcal S_NT_\epsilon.
\]
The first row of the equation gives exactly
\(T_\epsilon=\mathcal M_U b_\epsilon\). Hence both its thin and
repair components are actual polarized stresses, even before full
spectral closure. Free body coercivity and \(g^{-1}\le1\) give
\[
 \|g^{-1}T_\epsilon\|_Z
 \le P_\nu\|f\|_H.
\]
By \eqref{near:masked-free-bound},
\[
 \|g^{-1}b_\epsilon\|_H
 \le P_\nu e^{\eta_\delta/\nu}\|f\|_H.
\]
The remaining exponent can be made as small as required before viscosity
is chosen. These estimates do not assert that the intermediate masked
velocity is the final whole-space mode.
They apply to the prescribed body input. For an arbitrary equation
input of the form \(\mathcal M_Uv_{\rm in}\), the corresponding
identity is \(b=v_{\rm in}-\mathcal S_NT\), and its velocity bound
retains \(\|g^{-1}v_{\rm in}\|_H\). A small stress norm alone does
not control that velocity witness.

If a fixed smooth Gaussian affine force is used in both geometries,
write \(u_0=\mathcal M_0\mathcal R_0f\),
\(T_0=\mathsf D_0^{-1}u_0\). Suppose its \(H\) norm is a fixed
viscosity polynomial. The body version of the same output comparison,
with rate \(\epsilon^{1/14}\), and the coefficient/end-region
bounds above give
\(\|g^{-1}(u_\epsilon-u_0)\|_Z\le P_\nu e^{-c/\nu}\)
for some fixed \(c>0\). No gain factor is needed on the input here:
\(f\) is the prescribed unit force, not an amplified stress history.
The resolvent identity therefore yields
\[
 \|g^{-1}(T_\epsilon-T_0)\|_Z
 \le2Q\|g^{-1}(u_\epsilon-u_0)\|_Z
       +2Q^2\delta_N\|g^{-1}u_0\|_Z
 \le P_\nu e^{-c'/\nu}.
\]
In particular any old tensor
moment whose norm after gain multiplication is a viscosity polynomial
has the same limit on these two columns.
